\section*{October}

\section*{initial reading:}
$Key \ Vocabulary$ \\
Echo State Property:
\begin{itemize}
    \item reservoir's state is uniquely determined by the input history, not by the initial condition
    \item $x(t+1) = f(x(t), u(t)) \ and \ \norm{x(t)-x'(t)} \rightarrow 0 \ as \ t \rightarrow \infty$
\end{itemize}

Fading Memory:
\begin{itemize}
    \item reservoir's current state depends only weakly on inputs from the distant past
\end{itemize}

Input Separability:
\begin{itemize}
    \item ability to discriminate between different input sequences $\to$ adapt behavior to problem
\end{itemize}

Partial Trace:
\begin{itemize}
    \item tracing out the degrees of freedom of a subsystem
    \item $\rho_A = Tr_B(\rho_{AB})$ 
    \item $\ \rho_{AB}\in \mathcal{H}_A \otimes \mathcal{H}_B, \ \mathcal{H}_A, \mathcal{H}_B$ are the Hilbert spaces of subsystems A and B, respectively.
\end{itemize}

$RC \ Properties$
\begin{itemize}
    \item Echo State Property
    \item Fading Memory
    \item Input Separability
    \item fulfils all requirements $to$ universal
\end{itemize}

$Systems$ \\
Quantum Systems:
\begin{itemize}
    \item transverse fiel Ising model
\end{itemize}



$Online Protocol$ 
\begin{itemize}
    \item Input and feedback encoding of data and measurements into quantum reservoir
    \item  reservoir dynamics - system evolves under a definite quantum transformation
    \item continuous measurement and preparation of the reservoir for the next step without reset
\end{itemize}

\section*{19.10}
\textbf{Reading paper "Dissipation as a resource for QRC"}
\begin{itemize}
    \item engineered losses for each of the spin nodes $\to$ continuous dissipation
    \item dissipation crucial for echo state property and fading memory
    \item method: interactions with external environment under Marcovian conditions
    \begin{itemize}
        \item time evolution: Linbald master equation (GLKS) 
        \begin{flalign*}
            \dot{\rho} &= -i[H, \rho] + \sum_{j=1}^{N} \gamma_j
            \left(L_j \rho L_j^{\dagger}
            - \frac{1}{2}\{L_j^{\dagger}L_j, \rho\}\right) &&\\
            &= (\mathcal{U}+\mathcal{D})\rho
            \quad \text{with } \mathcal{U}\rho = -i[H, \rho]
            \text{ and } \mathcal{D}\rho = \text{Lindbladian term} &&
        \end{flalign*}

        \item $L_i$: Jump operator, $\gamma_i$: dissipation rate
        \item here $L_i = \sigma_i^- = \frac{1}{2}(\sigma_i^x - i\sigma_i^y)$
    \end{itemize}
    \item method: continuous dissipation (CD)
    \begin{itemize}
        \item new Hamiltonian $H'$ includes $h'(t) = h\cdot(s_k+1)$ $\to$ input signal dictates strength of of magnetic field to the x axis
        \item reservoir dynamics: $\dot{\rho} = (-i[H'(h_k'), \rho]+\mathcal{D}_L)\rho$
        \item uniform losses $\gamma_i = \gamma \forall i$
        \item local dissipator $\mathcal{D}_L = \gamma \sum_{i=1}^{N} \sigma_i^- \rho \sigma_i^+ - \frac{1}{2}\{\sigma_i^+\sigma_i^-, \rho\}$ \\ 
        $\to$ local because the chosen jump operators act on individual spins
        \item update rule: $\rho_{k+1} = e^{\Delta t(\mathcal{U}(s_{k+1})+\mathcal{D}_L)}\rho_k$
    \end{itemize}
    \item this generator is universal
    \item benchmark:
    \begin{itemize}
        \item tasks
        \begin{itemize}
            \item memory - short term memory, NARMA, parity check
            \item forecasting - Mackey-Glass
        \end{itemize}
        \item readout layer: linear combination of expectation values of Pauli strings of lenght 1 and 2
    \end{itemize}
\end{itemize}

\textbf{Reading paper "General theory of monitored QRC"}
\begin{itemize}
    \item develop general theory of monitored QRC based of indirect quantum measurements (weak, projective, partial, dissipative monitoring protocols)
    \item quantum monitoring: extracting information from quantum system through sequential measurements
    \item effect
    \begin{itemize}
        \item $E_m$ positive semidefinite 
        \item $\sum_m E_m = I$ (POVM) 
        \item $\operatorname{Pr}_\rho(m) = \operatorname{Tr}(E_m \rho)$
        \item $[E_m, X] = 0 \forall m \to$ useful for extracting information about observable $X$. 
    \end{itemize} 
    \item back-action: instead of strict projecttion: measurement-induced state disturbance $\to$ captured by Kraus operators $\{\Omega_m\}$ with $\Omega_m^\dagger \Omega_m = E_m$, where $E_m$ has the role of a "projector".
    \item measurement Map: $\mathcal{M}^{(m)}\rho = \frac{\Omega^\dagger \rho \Omega}{\operatorname{Pr}_\rho(m)}$
    \item measurement record $\vct{m}_{k} = (m_1, m_2, \ldots, m_k)$
    \item $\rho_k^{(m_1, m_2, \ldots, m_k)} = (\mathcal{M}^{(m_k)} \circ \Lambda_k) \rho_{k-1}^{(m_1, m_2, \ldots, m_{k-1})} $ with $\Lambda_k$: generig CTPT map describing the evolution of the system between measurements
    \item Gerald-Naimark theorem: $\mathcal{M}$ can always be implemented as an indirect measurement: Unitary coupling between system of interest (S) and quantum probe (P) followed by projective measurement on the probe. 
    \item Kraus representation of a CPTP map: \\ 
    $\mathcal{M}\rho = \sum_m \operatorname{Pr}_\rho (m) \cdot \mathcal{M}^{(m)}\rho = \sum_m \Omega_m \rho \Omega_m^\dagger $ \\
    $\to$ describes deterministic but irreversible dynamics
    \begin{enumerate}
        \item Depahsing measurement map: \\
        coupling unitary - $U_{PS} = CNOT_{S,P}(RY(\theta)_P \otimes I_S)$ 
        \item Amplitude damping measurement map \\
        coupling unitary - $U_{PS} = CNOT_{S,P}CRY(\theta)_{S,P}$ \\
        (controlled Y rotation)
        \item Partial measurement map: \\
        assign some of the units of a large system the role of the proble \\
        $\Omega_{\vct{b}_m}^{PM} = \Omega_{b_1} \otimes ... \otimes \Omega_{b_n} \otimes I^{\otimes (N-N_P)}$
    \end{enumerate}
    \item quantum channel $\mathcal{T}(s_k)$ such that $\rho_k = \mathcal{T}(s_k) \circ ... \circ \mathcal{T}(s_{1}) \rho_o$
    \item Learning task: weight matrix such that predicting sequence $\hat{\vct{y}}_k = \matr{W} \vct{x}_k$ minimizes loss to target sequence $\vct{y}_k$
    \item online protocol:
    \begin{itemize}
        \item the QRC map $\mathcal{T}^{(m)}(s) = (\mathcal{M}^{(m)}\circ \Lambda) (s)$
        \item feature vector: $x_{k,i} = \operatorname{Tr}[M_i\Lambda(s_k)\rho_{k-1}]$
        \item ESP and FMD given if both $\mathcal{M} \circ \Lambda$ is contractive 
        \item viable QRC map must be both stricty contractive and non-unital
        \item Measurement back-action con provide contractivity and non-unitality: \\
        Amplitude damping and partial projective measurements with reset work, Projective , dephasing and partial projective are unital 
    \end{itemize}
    \item uncertainty of observable mean values: projective measurements reach saturation faster, but weak measurements perform better in long term memory with sufficient shots
    \item time multiplexing: multiple virtual nodes with measurements, increases feature space $\to$ amplifies information-disturbance trade off
\end{itemize}


\section*{20.10}
\textbf{Call with Tobi}
\begin{itemize}
    \item discussion of the paper "General theory of monitored QRC"
    \item for online protocol to work $\to$ only one observable (pauli z)
    \item Programming Task: reproduce results of paper with online protocol (6 sping, half (3) control nodes)
\end{itemize}

\section*{21.10}

\textbf{Understanding tobi's qrcspin.py file:} \\
Idea:
n qubit with random Ising type Hamiltonian, including dissipation

Pipeline:
\begin{enumerate}
    \item encode data into hamiltonian via $\_generate\_field\_strength$ and build hamiltonian
    \item Evolve: $qt.mesolve$ integrates Linblad master eq., $L_i = \sqrt{\gamma} \cdot \sigma_i^- $ ($\gamma$: dissipation rate )
    \item Readout: expectation values, assumes 
    \item Optional diagnostics: negativity entaglement and l1-coherence of $\rho$
\end{enumerate}

Dynamics of the system: See paper "Dissipation as a resource for QRC" 

In simulation: Obtaining feature vector from $\rho$ does not collapse the state, as it is calculated exactly. 
Instead, the finite-shot mode mimics the statistical noise (sampling limited number of outcomes)

\textbf{Drafting plan for online protocol of the spin system:} \\
Reservoir
\begin{itemize}
    \item 6 spins, 3 probe spins, 3 reservoir spins
    \item System: $H'$ of "Dissipation as a resource for QRC" paper
    \item time evolution: $qt.mesolve$ with Linblad master equation
\end{itemize}

Online Protocol
\begin{itemize}
    \item Observables: Full set of Z strings
    \item Step: 1. Evolve, 2. Measure, 3. Reset probe spins (back-action), 4. Repeat
\end{itemize}


\section*{23.10}
Built the class $OnlineSpinQRC$ and  the test file $test\_online\_spin\_qrc.py$. \\
The class was built based on the class $SpinQRC$ but rewritten. It must still be adapted to 

\section*{24.10}
Summary of the Class $OnlineSpinQRC$:\\

\textbf{init:}\\
Important parameters:
\begin{itemize}
    \item model: choose "TiltedTFIM" (encoding in x direction)
    \item encoding method: chose "uniform" (for encoding method of "Dissipation as a resource for QRC" paper)
    \item $\theta := \pi \to$ projective measurement (back action is more generally implemented as a rotation)
    \item damping on: ?
    \item decay to: ? 
\end{itemize}

Initialization of basic operators:
\begin{itemize}
    \item $\sigma_i^x, \sigma_i^z, P0_i, P1_i, E01_i $
    \item combined with the following method for operator "X": \\
    \begin{lstlisting}
    self.X_ops = [self._build_operator(qt.X(), i) for i in range(n)]
    \end{lstlisting}
    $\to$ $[X_1 \otimes I_2 \otimes ... \otimes I_n, \ I_1 \otimes X_2 \otimes ... \otimes I_n, \ ... ]$
\end{itemize}

Random Number Generator:
\begin{itemize}
    \item needs to be explained
\end{itemize}

Shots:
\begin{itemize}
    \item needs to be explained
\end{itemize}

\textbf{run:}
running the class
\begin{enumerate}
    \item implement policy
    \item run the online protocol: change $\gamma_k, \theta_k$ for each step given the policy and perform an evolution step
    \item return the features
\end{enumerate}

\textbf{step:}
performing a single step of the online protocol
\begin{enumerate}
    \item evolve $\rho$ for given $\gamma_k$ and $s_k$
    \item obtain features $\to$ how does feature function work?
    \item measure and reset: apply the back action via the Kraus operators for given $\theta_k$
\end{enumerate}

OnlineSpinQRC.evolve():
\begin{itemize}
    \item contruct Hamiltonian
    \item Linblad evolution via $qt.mesolve$
\end{itemize}

OnlineSpinQRC.features():
\begin{itemize}
    \item how does this work?
\end{itemize}

OnlineSpinQRC.measureandreset():
\begin{itemize}
    \item back action on every probe: $\rho \to \sum K0_i \rho K0_i^\dagger = K0 \rho K0^\dagger + K1 \rho K1^\dagger$
    \item if ground state $\to$ let unchanged, if excited state $\to$ rotate back to ground state by angle $\theta_k \to U_{SA} = CRY(\theta_k){A,S} \cdot CNOT{S,A}$ (with additional purely logical auxillary qubit)
    \item $\to$ corresponding Kraus operators: $K0 = P0$ and $K1 = \cos(\theta)P1 + \sin(\theta)E01$
    \item literature?
\end{itemize}

tasks left for understanding:
\begin{itemize}
    \item understand dampingon $\to$ damping on depermines to which spins the lowering operator is applied to (meaning, which spins information can dissipate from)
    \item understanding decayto $\to$ to which state the lowering operator decays to $\ket{0}, \ket{1}$. Convention: $\sigma_z \ket{0} = +\ket{0}$ and $\sigma_z \ket{1} = -\ket{1}$ and thus $\ket{0}$ is the higher energy state and $\ket{1}$ is the lower energy state \\$\to$ default: decay to = 1.    \item understand the feature function and the shot system
    \item look for literature for the back action
\end{itemize}

new tasks
\begin{itemize}
    \item unit tests
    \item test, if it implements the online protocol and the reservoir correctly
    \item prepare a read me for tobi
\end{itemize}

\section*{27.10}
Meeting with tobi - discussed the paper, delete unesccessary functionalities from the class

\section*{29.10}
\begin{itemize}
    \item added literature and thesis to the git folder
    \item removed unused encoding methods and class variable encoding strength
\end{itemize}